% META: {"id": "TNSI-ALGO-COURS-03", "chapitre": "TNSI-ALGORITHMIQUE", "type_objet": "cours", "capacites_codes": ["C3"], "sources_inspiration": [], "mode_creation": "ex_nihilo", "status": "approved", "capacites": ["TNSI-ALGORITHMIQUE-C3"]}
\section{Parcours d'un graphe}

% ============================================================
% STRATE 1 — BFS
% ============================================================

On représente un graphe par un dictionnaire associant à chaque sommet la liste de ses successeurs (voir le chapitre \textit{Structures de données}).

\definition[1]{
Le parcours en \textbf{largeur d'abord} (\textit{Breadth-First Search}, BFS) d'un graphe visite les sommets par distance croissante au sommet de départ, en utilisant une \textbf{file}. Contrairement à un arbre, un graphe peut contenir des \textbf{cycles} : il faut donc mémoriser les sommets déjà visités pour ne pas boucler indéfiniment.
}

\textbf{Exemple.}
\begin{python}
def bfs(graphe, depart):
    visites = [depart]
    file = [depart]
    while file:
        sommet = file.pop(0)
        for voisin in graphe[sommet]:
            if voisin not in visites:
                visites.append(voisin)
                file.append(voisin)
    return visites
\end{python}

% BEGIN-VERIFY
% def bfs(graphe, depart):
%     visites = [depart]
%     file = [depart]
%     while file:
%         sommet = file.pop(0)
%         for voisin in graphe[sommet]:
%             if voisin not in visites:
%                 visites.append(voisin)
%                 file.append(voisin)
%     return visites
%
% graphe = {0: [1, 2], 1: [0, 3], 2: [0, 3], 3: [1, 2]}
% assert bfs(graphe, 0) == [0, 1, 2, 3]
% END-VERIFY

\margeAppui{
Sans la liste \lstinline{visites}, un graphe contenant un cycle (comme le carré $0-1-3-2-0$ de l'exemple) ferait boucler l'algorithme indéfiniment : le sommet $0$ serait redécouvert depuis $1$ ou $2$, encore et encore.
}

% ============================================================
% STRATE 2 — DFS
% ============================================================

\definition[2]{
Le parcours en \textbf{profondeur d'abord} (\textit{Depth-First Search}, DFS) explore chaque branche aussi loin que possible avant de revenir en arrière (\textit{backtrack}). Il s'implémente naturellement par \textbf{récursion} (la pile d'appels joue le rôle de la pile), ou explicitement avec une \textbf{pile}.
}

\textbf{Exemple.} Version récursive :
\begin{python}
def dfs(graphe, sommet, visites=None):
    if visites is None:
        visites = []
    visites.append(sommet)
    for voisin in graphe[sommet]:
        if voisin not in visites:
            dfs(graphe, voisin, visites)
    return visites
\end{python}

% BEGIN-VERIFY
% def dfs(graphe, sommet, visites=None):
%     if visites is None:
%         visites = []
%     visites.append(sommet)
%     for voisin in graphe[sommet]:
%         if voisin not in visites:
%             dfs(graphe, voisin, visites)
%     return visites
%
% graphe = {0: [1, 2], 1: [0, 3], 2: [0, 3], 3: [1, 2]}
% assert dfs(graphe, 0) == [0, 1, 3, 2]
% END-VERIFY

\erreurFrequente{
\textbf{Argument par défaut mutable.} Écrire \lstinline{def dfs(graphe, sommet, visites=[]):} au lieu de \lstinline{visites=None} est un piège classique en Python : la liste par défaut est créée \textbf{une seule fois}, à la définition de la fonction, et \textbf{partagée} entre tous les appels qui ne fournissent pas explicitement \lstinline{visites}. Un deuxième appel \lstinline{dfs(graphe, 0)} repartirait alors avec les sommets déjà visités lors du premier appel.
}

\propriete[BFS et DFS : deux usages différents]{
Le BFS trouve le \textbf{plus court chemin} en nombre d'arêtes dans un graphe non pondéré (il visite par distance croissante). Le DFS est souvent plus simple à écrire et suffit pour explorer entièrement un graphe ou détecter sa connexité, mais ne garantit pas de trouver le chemin le plus court.
}
